\scanpage{283}

\unsure{suppose} moreover \deleted{\unsure{that} \illegible{4}}\inserted{(2°) \deleted{\illegible{1}} $1,\rho,\sigma$ \unsure{are} linearly} independent in each fiber i.e. $1,\rho^{-1},N_1$ linearly independent in each fiber i.e. $1,\sigma^{-1},N_1$ linearly independent in each fiber, \deleted{\illegible{4}}\inserted{and (3°) $\operatorname{Tr}N_1\rho$ invertible.}

\margin{\illegible{2} $1,\rho^{-1},\varepsilon_1(=1+N_1)$ \illegible{1}. \illegible{2} $1,\sigma^{-1},\varepsilon_1(=1+N_1)$ \illegible{1}.}

\margin{\uline{N.B.} 2°) is a consequence of 3°), which implies \unsure{primo}, \illegible{3} \unsure{suitable} \unsure{basis}$^*$
\[
\varepsilon_1=\begin{pmatrix}1&0\\1&1\end{pmatrix},\qquad
\rho=\begin{pmatrix}0&1\\a&b\end{pmatrix}\quad (a\text{ invertible}),
\]
\[
\deleted{\illegible{1}}=N_1\rho=\unsure{\cdots}=\begin{pmatrix}0&0\\0&1\end{pmatrix},
\qquad\text{thus }\sigma=\begin{pmatrix}0&-1\\-a&1-b\end{pmatrix}.
\]
$^*$ and on \unsure{normalizing} so as to have $\operatorname{Tr}N_1\rho=1$.}

(N.B. we have $N_1\rho=-N_1\sigma=\rho+\sigma$). This being so, we will have
\[
\operatorname{int}(\rho^{-1})\varepsilon_1=\operatorname{int}(\sigma^{-1})\varepsilon_1,\quad\text{i.e. }\varepsilon_0.
\]
(thus $\operatorname{int}(\rho)\varepsilon_0=\operatorname{int}(\sigma)\varepsilon_0=\varepsilon_1$)

\deleted{Let $B_0$ be the normalizer of \illegible{1}\inserted{$L_0=\unsure{\text{subgroup}}$} let $B'_0\subset B_0$ be the subgroup scheme defined by (74), so that we have}
\[
\varepsilon_0=1+N_0.
\]
\deleted{\unsure{set}} \inserted{set}
\[
\varepsilon_0(\lambda)=1+\lambda N_0,\qquad\varepsilon_1(\lambda)=1+\lambda N_1,
\]
(additive 1-parameter subgroups)
and let $B_0$ be the \inserted{(Borel)} normalizer \inserted{in $G(\mathcal M)$} of the
\[
\begin{aligned}
L_0&=\text{subgroup scheme }\varepsilon_0(\mathrm G_a)\\
L_1&=\text{\rule{6.3cm}{0.2pt}}\quad\varepsilon_1(\mathrm G_a),
\end{aligned}
\]
(thus $\rho(L_0)=\sigma(L_0)=L_1$),
let $B'_0$ be the subgroup scheme of $B_0$ defined \inserted{in terms of $\varepsilon_0$,} by (74), so that $B_0$ fits into an exact sequence
\[
1\longrightarrow L_0\hookrightarrow B_0\xrightarrow{\det}\mathrm G_m\longrightarrow1,
\]
and that $\underline u\in\Gamma B_0$ implies
\[
\underline u(N_0)=\det(\underline u)N_0,\qquad\text{i.e.}
\]
\[
\underline u(\varepsilon_0(\lambda))=\varepsilon_0(\det(\underline u)\lambda).
\]

This being so, consider the morphism
\[
\begin{aligned}
B'_0&\longrightarrow\deleted{\unsure{\mathcal X}}G(\mathcal M)\times G(\mathcal M)\times\unsure{\mathbb E^1|\mathbb E'}\\
\underline u&\longmapsto([\underline u,\rho],[\underline u,\sigma],\det(\underline u)).
\end{aligned}
\]
\mspage{282}

\scanpage{284}
\mspage{443}

This takes its values in $\mathcal X_{\rho,\sigma}$, the scheme of solutions \deleted{\unsure{of the}} simultaneous equations (33) to (35), and induces an isomorphism
\[
\tag{79}B'_0\unsure{\simeq|\xrightarrow{\sim}}\mathcal X_{\rho,\sigma}.
\]

\deleted{\unsure{Corollary.}} Let us recall \unsure{also|thus} for the record

\uline{Corollary} The projection
\[
\begin{aligned}
\mathcal X_\rho&\longrightarrow G(\mathcal M)\\
(\zeta,\eta,\mu)&\longmapsto\zeta
\end{aligned}
\]
is a monomorphism, and the image consists of the $\zeta$ satisfying
\[
\det\zeta=1,\quad \operatorname{Tr}\zeta\rho=\operatorname{Tr}\rho,\quad \operatorname{Tr}\zeta\underbrace{(\rho+\sigma)}_{\inserted{\substack{\Vert\\-N_1\sigma=\unsure{+}N_1\rho=\rho N_0=-\sigma N_0}}}\text{ invertible}
\]

Likewise the projection
\[
\begin{aligned}
\mathcal X_\rho&\longrightarrow G(\mathcal M)\\
(\zeta,\eta,\mu)&\longmapsto\eta
\end{aligned}
\]
is a monomorphism, and the image consists of the $\eta$ satisfying
\[
\det\eta=1,\quad\operatorname{Tr}\eta\sigma=\operatorname{Tr}\sigma,\quad\operatorname{Tr}\eta(\rho+\sigma)\text{ invertible}
\]

Finally, we will have
\deleted{\illegible{1}$\det(\underline u)=$}
\[
\tag{80}\deleted{\unsure{\zeta}}\operatorname{Tr}\zeta(\rho+\sigma)=\operatorname{Tr}\eta(\rho+\sigma)=1/\mu.
\]

\uline{Scholium} By transport of structure, we find a \unsure{group|groups} structure on $\mathcal X_{\rho,\sigma}$.
\mspage{283}

\scanpage{285}

Making it explicit, we find
\[
\begin{array}{rcl}
\text{if }\underline u&\text{corresponds to}&\zeta,\eta,\mu\\
\underline u'&\text{---}&\zeta',\eta',\unsure{\mu'|\mu''}\\
\underline u''=\underline u'\underline u&\text{---}&\zeta'',\eta'',\mu''
\end{array}
\]
we will have
\[
\tag{\unsure{81|21}}\mu''=\mu\mu',\quad\zeta''=u'(\zeta)\zeta',\quad\eta''=u'(\eta)\eta'
\]
where we set, as usual
\[
u'(\unsure{x|X})=u'\unsure{x|X}u'^{-1},
\]
\deleted{\unsure{for}}
for two elements $x,u'$ of a group.

\uline{$\overline{\mathrm{III}}$}) Study of the equation (Continued)

\[
\tag{82}\alpha\rho(\alpha)^{-1}=\beta\sigma(\beta)^{-1}\varepsilon_1(\mu-1)
\]
\[
(\alpha,\beta\in\Gamma G(\mathcal M),\quad\mu\in\Gamma\unsure{\mathbb E_S^1|\mathcal O_S^*})
\]
i.e.
\[
\tag{82bis}[\alpha,\rho]=[\beta,\sigma]\varepsilon_1(\mu-1).
\]

Let $\mathcal G_{\rho,\sigma}$ be the scheme of solutions, we thus have a morphism
\[
\tag{83}\begin{aligned}
\mathcal G_{\rho,\sigma}&\longrightarrow\mathcal X_{\rho,\sigma}\\
(\alpha,\beta,\unsure{\mu})&\longmapsto(\alpha\rho(\alpha)^{-1},\beta\deleted{\unsure{\rho}}\inserted{\sigma}(\beta)^{-1},\mu)
\end{aligned}
\]
and it follows from the preceding theorem that this is \deleted{\illegible{1}} an epimorphism, more precisely,
\mspage{284}

\scanpage{286}
\mspage{444}

We have a canonical section
\[
\tag{84}
\diagram{\begin{tikzcd}[ampersand replacement=\&,column sep=large,row sep=small]
\mathcal G_{\rho,\sigma}\arrow[r]\&\mathcal X_{\rho,\sigma}\&([\underline u,\rho],[\underline u,\sigma],\det\underline u)\\
\&B'_0\arrow[ul]\arrow[u,"\wr"']\&\underline u\arrow[u,mapsto]\\
(\underline u,\underline u,\det\underline u)\&\underline u\arrow[l,mapsto]\&
\end{tikzcd}}
\]

This section takes its values in the ``diagonal'', \deleted{\unsure{form}} which corresponds to the pairs $(\underline u,\mu)$ \inserted{$(\underline u\in\Gamma\mathcal M^*,\ \mu\in\Gamma\unsure{\mathbb E^1|\mathbb E'})$} such that
\[
\tag{85}[\underline u,\rho]=[\underline u,\sigma]\varepsilon_1(\mu-1),
\]
which we \unsure{will try|try} to determine as well.

Combining the results of (I, and II) we see that the general solution of (82) will be \unsure{formed|provided} by the,
\[
(\underline u\underline a,\underline u\underline b,\mu=\det\underline u)
\]
where $\underline u$ is a section of $B'_0$, $\underline a$ a \deleted{local} section of $T_\rho$, $\underline b$ a section of $T_\sigma$.

To be precise

\uline{Theorem.} Consider
\[
\tag{86}\begin{aligned}
B'_0\times T_\rho\times T_\sigma&\longrightarrow G\times G\times\unsure{\mathbb E^1|\mathbb E'}\\
(\underline u,\underline a,\underline b)&\longmapsto(\underline u\underline a^{-1},\underline u\underline b^{-1},\det\underline u)
\end{aligned}
\]
\mspage{285}

\scanpage{287}

\unsure{It} takes its values in $\mathcal G_{\rho,\sigma}$, and is a monomorphism --- hence induces an \uline{isomorphism}
\[
\tag{87}B'_0\times T_\rho\times T_\sigma\xrightarrow{\sim}\mathcal G_{\rho,\sigma}
\]

More precisely, let \deleted{\unsure{locally}} $(\alpha,\beta,\mu)$ be a section of $\mathcal G_{\rho,\sigma}$, then there exist \unsure{then} \inserted{\unsure{some}} unique sections \deleted{$\underline a$ of} $\underline a$ of $T_\rho$, $\underline b$ of $T_\sigma$, such that we have
\deleted{$\alpha\underline a=\beta\underline b\ (\text{i.e. }\underline u)$}
\[
\tag{88}\alpha\underline a,\beta\underline b\in\Gamma B'_0
\]
\deleted{$\beta\underline b$ \illegible{2}}
and we will then necessarily have
\[
\alpha\underline a=\beta\underline b\ (\overset{\mathrm{df}}{=}\underline u)
\]
whence of course
\[
([\alpha,\rho],[\beta,\rho])=([\underline u,\rho],[\underline u,\deleted{\unsure{\rho}}\inserted{\sigma}])
\]
and the unique $\mu$ satisfying (82) is therefore $\det\underline u$ \deleted{\illegible{1}} ; thus
\[
\tag{\unsure{89}}(\alpha,\beta,\mu)=(\underline u\underline a^{-1},\underline u\underline b^{-1},\det\underline u)\quad\text{(uniquely)}
\]
\deleted{\unsure{Proof.}\ \unsure{It}}

Proof. It remains only to make explicit the uniqueness of $\alpha,\beta$ satisfying (88), which is a consequence of
\[
\tag{90}T_\rho\cap B'_0=T_\sigma\cap B'_0=\{1\}
\]
which we had forgotten to make explicit in (II).
\mspage{286}

\scanpage{288}
\mspage{445}

In terms of the \inserted{sole} datum of $\alpha$, we will therefore have an $\underline a$ uniquely determined by the condition $\alpha\underline a\in\Gamma B'_0$, and we will have
\[
(91)\qquad \underline a=c_0\rho+d_0
\qquad
\begin{gathered}
c_0,d_0\in\Gamma(\mathcal O_S)\\
\text{determined by }\alpha
\end{gathered}
\]
\unsure{likewise} we will have
\[
(92)\qquad \underline b=t_0\sigma+u_0
\qquad
\begin{gathered}
t_0,u_0\in\Gamma(\mathcal O_S)\\
\text{determined by }\beta
\end{gathered}
\]
in such a way therefore that we will have
\[
(93)\qquad \alpha\unsure{\underline a}=\beta\underline b
\qquad (\overset{\mathrm{dfn}}{=}\underline u)
\]

In the \unsure{numerical} case
\[
\rho=\begin{pmatrix}0&1\\-1&1\end{pmatrix},\qquad
\sigma=\begin{pmatrix}0&-1\\1&0\end{pmatrix},
\]
setting
\[
\alpha=\begin{pmatrix}a&b\\c&d\end{pmatrix},\qquad
\beta=\begin{pmatrix}r&s\\t&u\end{pmatrix}
\]
we will have (cf. page\ $398'$ (76)(78)) unless I am mistaken
\[
(94)\qquad \underline a=\frac1{\delta D}(c\rho+d)
\]
\inserted{\[
\left(=\frac1{c^2+cd+d^2}(c\rho+d)\right)
\]}
where $\delta D=c^2+cd+d^2$, --- $D$ being the \unsure{corner} coefficient of
\[
\zeta=\alpha\rho(\alpha)^{-1}=\begin{pmatrix}A&B\\C&D\end{pmatrix}
\]
\[
(95)\qquad \underline b=\frac1{\delta' U}(t\sigma+u)
\]
\inserted{\[
=\frac1{t^2+u^2}(t\sigma+u)
\]}
where $\delta' U=t^2+u^2$, --- $U$ the corner coefficient of
\[
\eta=\beta\sigma(\beta)^{-1}=\begin{pmatrix}R&S\\T&\unsure{U|u}\end{pmatrix},
\]
\margin{$\delta=\det\alpha=ad-bc$\par
$\delta'=\det\beta=ru-st$}
\margin{We will have to come back to these formulas in the general case}

the only trouble, in making the link with the ``brutal'' computations of loc. cit, is in the factor $\tau_\infty^\alpha$ on page\ $398'$ (78), where $\alpha=\unsure{\alpha_2}(\mu)\in\mathbb Z/2.\mathbb Z$, in the case where $\alpha=1$ ---

\mspage{287}

\scanpage{289}

at first sight there seems to be an incompatibility between the (simpler) theory developed here, and this factor $\tau_\infty^\alpha$ --- which \unsure{I} must think has very good reasons for being there ! And there it is --- in loc. cit. we wrote the equation \unsure{of the} \unsure{loops} with the possibility of a \uuline{sign}
\[
\zeta=\varepsilon\eta
\underbrace{\ell_1^\nu}_{\inserted{\substack{\varepsilon_1(\mu-1)\ \text{with the}\\\text{current notation}}}},
\qquad\varepsilon\in\{\pm1\},
\]
\margin{here $\zeta$ and $\eta$ are \deleted{\illegible{1}} defined as $\alpha\rho(\alpha)^{-1}$, $\beta\sigma(\beta)^{-1}$}
this is thus a \uline{generalization} \unsure{relative} to the equation set down at the beginning, where we assumed $\varepsilon=+1$.
We will therefore have to \unsure{later} take up the preceding theory again, in the more general context of an equation
\[
(\unsure{95|96})\qquad
\underset{\inserted{\substack{\Vert\\\alpha\rho(\alpha)^{-1}}}}{\zeta}
=\varepsilon.\underset{\inserted{\substack{\Vert\\\beta\sigma(\beta)^{-1}}}}{\eta}\varepsilon_1(\mu-1)
\qquad\varepsilon\ \text{section of }\mathcal O_S
\]
where we will have to have (taking the determinants of the two sides) \deleted{$\varepsilon^2=1$ \illegible{1}}
\[
\deleted{\unsure{\det}}\qquad\varepsilon^2=1
\]
--- but there is probably no reason to suppose that we have $\varepsilon\in\{\pm1\}$. But I will first continue the calculations in the present case $\varepsilon=+1$.

\mspage{288}

\scanpage{290}
\mspage{446}

\uline{$\overline{\mathrm{IV}}$}\qquad\uline{A canonical pinning}

Let us recall that giving \inserted{(over a base scheme $S$)} a form $G$ of $\mathrm{Gl}(2)_S$, or $SG$ of $\mathrm{Sl}(2)_S$, \inserted{or $PG$ of $\mathrm{GP}(1)_S$,} or again an \unsure{Azumaya|Azumaya} algebra $\mathcal M$ (form of $M_2(\mathcal O_S)$), amounts to giving a bundle of \unsure{projective} lines $\unsure{\mathcal P|P}$ over $S$.
$\unsure{\mathcal P|P}$ can be interpreted as the scheme of Borel subgroups \inserted{$B$} of $G$, or \inserted{$SB$} of $SG$, \inserted{or $PB$ of $PG$,} or as the scheme of ``one-parameter \inserted{additive} ``Borelian'' subgroups'' (i.e. identical with the unipotent part of a suitable Borel) \deleted{\illegible{1}}\inserted{$L$} of $G$, or of $SG$, or of $PG$ (N.B. The canonical morphisms $SG\longrightarrow G\longrightarrow PG$ induce bijections between \inserted{the \unsure{sets} of} Borel subgroups of the three groups in question, and also between the \unsure{sets} of ``Borelian'' 1-parameter subgroups of these groups).
Likewise, $\unsure{\mathcal P|P}$ can be interpreted in terms of the sections $N$ \deleted{\illegible{1}} of $\mathcal M$ which are nonzero on \unsure{every} fiber and such that \deleted{$N^2=0$},
\inserted{\[
(1)\qquad\det N=0,\quad\operatorname{Tr}N=0\qquad\text{hence }N^2=0
\]}
(the ``strictly \unsure{nilpotent}'' \inserted{and regular (nonzero on each fiber)} sections), to every such section \unsure{will correspond|corresponds} a \inserted{1-parameter} Borelian subgroup of $SG$, as the image of $\unsure{\mathbf G_a|\mathrm G_a}$ under
\[
(2)\qquad\lambda\longmapsto\inserted{\unsure{e_N}(\lambda)=}\exp(\lambda N)=1+\lambda N.
\]
But it is obvious that if $N$ and $N'$ are two such

\mspage{289}

\scanpage{291}

sections which are homothetic, i.e. $N'=aN$ with $a\in\Gamma(S,\mathcal O_S)$, \inserted{(in fact, we will have $a\in\Gamma(S,\mathcal O_S^*)$)} then they define the \unsure{same} 1-parameter subgroup, --- so the latter in fact correspond \deleted{to the}\inserted{to} sections of the projective bundle (of relative dimension 3) $\check P(\mathcal M)$, or if we take $\operatorname{Tr}N=0$ into account, \deleted{\illegible{1}} to sections of the projective bundle (of relative dimension 2 --- thus a bundle of projective planes) $\check P(\unsure{{}_1\mathcal M|1\mathcal M})$ (where $\unsure{{}_1\mathcal M|1\mathcal M}\subset\mathcal M$ is defined precisely by the condition $\operatorname{Tr}N=0$) --- namely the sections of the \deleted{\unsure{quadric}} ``conic'' defined by the equation
\[
\det N=0.
\]

The $N$'s moreover also correspond to the operators
\[
e=e_N=\deleted{\unsure{1+N}}e_N(1)=1+N
\]
satisfying the conditions
\[
(3)\qquad\det e=1,\quad\operatorname{Tr}e=2
\]
which is indeed equivalent, on setting $e=1+N$, i.e. on defining $N$ by
\[
N=e-1
\]
to the conditions (1) [for $\operatorname{Tr}$ and $\det$ \unsure{of two|of the} two operators $N$ and $e$ related by $e=\lambda+N$ mutually determine \deleted{\unsure{one another}} by
\[
\left\{
\begin{aligned}
\operatorname{Tr}e&=2\operatorname{Tr}\lambda+\operatorname{Tr}N,&
\det e&=\lambda^2+\lambda\operatorname{Tr}N+\det N\\
\operatorname{Tr}N&=-2\operatorname{Tr}\lambda+\operatorname{Tr}e,&
\det N&=\lambda^2-\lambda\operatorname{Tr}e+\det e
\end{aligned}
\right.
\quad .
\]

\mspage{290}

\scanpage{292}
\mspage{447}

The \deleted{operators} sections of $G$ (in fact of $SG$) satisfying (3), \inserted{and which are not $=1$ on any fiber,} are called the \deleted{\unsure{strict|strictly}} \inserted{regular} unipotent sections --- two such sections \inserted{$e,e'$} are considered ``equivalent'' if we have
\margin{Steinberg's terminology}
\[
\begin{aligned}
(4)\qquad e'&=e^\lambda\deleted{\unsure{=e}}\inserted{\overset{\unsure{\mathrm{def}|\mathrm{dfn}}}{=}e(\lambda)}
&& (\lambda\in\Gamma\unsure{\mathbf G_a|\mathrm G_a}=\Gamma\mathcal O_S)\\
&\overset{\mathrm{def}}{=}(1+\lambda N)
&&\begin{gathered}\deleted{\text{we will have }\unsure{\text{always}}}\\\text{we will have }\unsure{\text{necessarily}^{\mathrm t}}\quad\lambda\in\Gamma\unsure{\mathbf G_m|\mathrm G_m}\end{gathered}\\
&=1+\lambda(e-1)=(1-\lambda)+\lambda e.
\end{aligned}
\]

We thus see that the regular (strictly) nilpotent sections of $\mathcal M$ correspond \inserted{1-1} to the regular unipotent sections, and the equivalence relation of homothety becomes that of equivalence by $e\longmapsto e^\lambda$ --- and the equivalence classes of either correspond to the Borelian additive subgroups of $G$ (or $SG$, or $PG$ via $SG\longrightarrow PG$) --- or again $\longrightarrow$ (taking the normalizers of these in $G$, $SG$ or $PG$) to the Borels of these groups, i.e. to the sections of $\unsure{\mathcal P|P}$ over $S$.

Suppose such a section $\unsure{x_\infty}$ of $\unsure{\mathcal P|P}$ is given, this thus allows us to \deleted{see} \unsure{consider} $\unsure{\mathcal P|P}$ as obtained from an affine \inserted{line} bundle
\deleted{\illegible{3}}
\[
(5)\qquad E\quad(=\unsure{\mathcal P|P}\setminus\{\unsure{x_\infty}\}),
\]
\deleted{\unsure{A vector bundle} \illegible{4} vector of translations \illegible{3} \unsure{invertible} $D$}

\mspage{291}

\scanpage{293}

and $x_\infty$ is also given by a Borelian 1-parameter subgroup
\[
L\subset SG\subset G \tag{6}
\]
(which, under $SG\longrightarrow PG$, again gives a Borelian \unsure{additive} subgroup $L\hookrightarrow PG$ in $PG\ldots$). Locally, $L$ will be generated by a regular unipotent section
\[
e=1+N\qquad (N\text{ regular \unsure{nilpotent}}) \tag{7}
\]
as the \unsure{subgroup} of the
\[
e^\lambda=\unsure{\underline e(\lambda)|e_{\lambda}(\lambda)}=e_N(\lambda)=1+\lambda N, \tag{8}
\]
but as we have said, $e,N$ are not determined \unsure{intrinsically} in terms of $L\subset SG$. More precisely, $L$ is canonically endowed with a vector bundle structure (\unsure{preserved} by its normalizer in $G$), \deleted{\illegible{1}} \unsure{to|:} wit precisely that obtained by transport of structure from that of $G_a$ under
\[
\underline e:G_a\xrightarrow{\sim}L\qquad
\underline e(\lambda)=e^\lambda=1+\lambda N,
\]
and \deleted{the choice of} $L$ is thus identified with the vector bundle \unsure{canonically|canonical} associated with the line bundle $\mathcal N\subset\unsure{{}_1M}$ \deleted{\unsure{sub}} defined by the $N$'s
\[
L\simeq\check V(\mathcal N), \tag{9}
\]
\mspage{292}

\scanpage{294}
\mspage{448}

the choice of $e$, or of $N$, corresponds simply to that of a \uuline{basis} \inserted{$N$} of $\mathcal N$, or of a section \inserted{$e$} of $L^*$ \deleted{\unsure{everywhere} $\unsure{\ne}$} \inserted{i.e. of a section $e$ of $L$} which is $\ne1$ on \unsure{every} fiber.

\margin{Let us note that $L$ is also \unsure{canonically} \unsure{isomorphic} to the vector bundle of translations of the affine line bundle $\unsure{\mathbb P|P}\setminus\{\unsure{x|x_\infty}\}$.}

Let us return to equation (32), and note that $\rho$ and $\sigma$ enter into it only modulo \unsure{transformations} $\rho\mapsto\lambda\rho$, $\sigma\mapsto\lambda'\sigma$ \inserted{($\lambda,\lambda'\in\Gamma G_m$)} (which do not change the commutators $[\alpha,\rho]$ and $[\beta,\sigma]$); \deleted{on the other hand, \unsure{the data}} \inserted{We can express this by saying that $\rho,\sigma$ enter only through} \deleted{in a Borelian additive subgroup $L_1$ of $G$} their images in $\Gamma(PG)$, and moreover we have
\[
\begin{aligned}
PG&\xrightarrow{\sim}\underline{\mathrm{Aut}}_{\unsure{Sgr|S\text{-gr}}}(G)
\simeq\underline{\mathrm{Aut}}_{S\text{-gr}}(SG)
\simeq\underline{\mathrm{Aut}}_{\unsure{Sgr|Sg}}(PG)\\
&\simeq\underline{\mathrm{Aut}}_{S\text{-Alg}}(M)
\simeq\underline{\mathrm{Aut}}_S(\unsure{\mathbb P|P})
\end{aligned} \tag{\unsure{10|40}}
\]
and our remark amounts to saying that $\rho$ and $\sigma$ enter only through their action on $G$ (or on $SG$), which is also clear from the formulas
\[
[\alpha,\rho]=\alpha\rho(\alpha)^{-1},\qquad
[\beta,\sigma]=\beta\sigma(\beta)^{-1}.
\]

On the other hand, the factor $\varepsilon_1(\mu-1)$ of (32) can be regarded as a section of a \deleted{\unsure{false bundle}} additive subgroup $L_1$ of $G$ (in fact, of $SG$), which we suppose of
\mspage{293}

\scanpage{295}

course Borelian --- equation (32) is then written in the form (where $\mu$ no longer enters, only $\alpha$ and $\beta$)
\[
[\beta,\sigma]^{-1}[\alpha,\rho]\in\Gamma L_1 \tag{11}
\]
or
\[
[\deleted{\unsure{\beta}}\inserted{\alpha},\rho]=[\beta,\sigma]\cdot\gamma
\qquad\text{with }\gamma\in\Gamma L_1. \tag{12}
\]
In the form (12), we can regard it as an equation relating $\alpha,\beta,\gamma$.

Let us consider the image \deleted{$\unsure{\widetilde L_1|\widetilde L_0}$}
\[
\widetilde L_1\subset PG\deleted{\unsure{\ell}} \tag{13}
\]
of $L_1$ under $G\longrightarrow PG$, the relation (40 bis) is then written in the form
\[
\widetilde\varepsilon_1\overset{\unsure{\mathrm{dfn}}}{=}
\widetilde\rho\widetilde\sigma^{-1}\in\Gamma\widetilde L_1
\qquad(\xleftarrow[\unsure{\mathrm{can}}]{\sim}\Gamma L_1). \tag{14}
\]
\deleted{I say that \illegible{1} the relation}
This section therefore comes from a section \uline{well determined} \deleted{\unsure{and}} in terms of $\widetilde\rho,\widetilde\sigma$
\[
\varepsilon_1\in\Gamma L_1, \tag{15}
\]
\inserted{described in terms of $\rho,\sigma$ directly} by the condition that we have a relation of the form
\[
\varepsilon_1=-\lambda_0\rho\sigma^{-1}
\qquad(\lambda_0\in\Gamma G_m) \tag*{(16) = (40 bis)}
\]
\inserted{\uline{for convenience}} --- (because of the application to the special case that interests us most \ldots).

\margin{\uline{N.B.} $\lambda_0$ is uniquely determined by $\rho,\sigma$ if $2$ invertible on $S$; more \unsure{general\textsuperscript{ly}}, if $2$ regular \deleted{\unsure{section}}\inserted{\unsure{element}} of $\Gamma(S,\mathcal O_S)$. (\uline{But} \uline{in characteristic 2 $\unsure{\lambda_0}$ \unsure{is not}} \uline{\unsure{determined} \unsure{except}} \uline{modulo $\mu_2$} !!!).}
\mspage{294}

\scanpage{296}
\mspage{449}

Let us return to the general situation of a form $G$ of $\unsure{\mathrm{Gl}(2)_S|\mathrm{GL}(2)_S}$, and suppose that it can be written in the form
\[
G=\underline{\mathrm{Aut}}_{\mathcal O_S\text{-Mod}}(F), \tag{17}
\]
where $F$ is an $\mathcal O_S$-Module \unsure{local\textsuperscript{ly}} free of rank 2,

This is the case as soon as we have a section $x$ of $\unsure{\mathbb P|P}$ over $S$, i.e. a Borel additive subgroup $L$ of $G$, or of $SG$, (in fact we can then take for $F$ the extension of $\mathcal O_S$ by $\mathcal N$ which determines the affine-line bundle $\unsure{\mathbb P|P}\setminus\{x\}$). But \unsure{the|the} $\unsure{\mathbb P|P}$ is realized as \deleted{\illegible{1}}
\[
\unsure{\mathbb P|P}=\unsure{\mathbf P|P}(F)\simeq\unsure{\check{\mathbf P}|\check P}(F) \tag{18}
\]
the projective bundle associated with $F$, \deleted{\unsure{Thus}} \inserted{\unsure{i.e.}} the sections of $\unsure{\mathbb P|P}$ are identified with the ``lines'' $D\subset F$.

We will have
\[
M=\underline{\mathrm{End}}_{\mathcal O_S\text{-mod}}(F)\deleted{\unsure{\simeq}}\simeq\check F\otimes_{\mathcal O_S}F \tag{19}
\]
and the sub-\deleted{bundle}\inserted{Module} \deleted{vector}\inserted{$\mathcal N$} associated with the section $x_D$ of $\unsure{\mathbb P|P}$ \deleted{\unsure{is}} is none other than \deleted{\illegible{1}}
\[
\mathcal N=D^\circ\otimes D
\qquad\bigl(D^\circ\simeq(F/D)^\vee\subset\check F,\text{ the orthogonal to }D\text{ in }\check F\bigr) \tag{20}
\]
its sections are the $u:F\longrightarrow F$ such that
\[
\deleted{u(F)\subset D,\quad u(D)=0,}
\]
\[
u(F)\subset D,\quad u(D)=\{0\}. \tag{21}
\]
\mspage{295}

\scanpage{297}

\uline{Proposition} \deleted{Under the preceding conditions} \inserted{Let $G$ be a form of $\unsure{\mathrm{Gl}(2)_S|\mathrm{GL}(2)_S}$, $\rho,\sigma$ two sections of $G$ such that $\widetilde\rho\widetilde\sigma^{-1}=\widetilde\varepsilon_1$ is a \inserted{\illegible{1}} \unsure{unipotent}-\deleted{\unsure{regular}} section of $PG$, \unsure{linked} \unsure{by the} relation (16). Then} the following conditions are equivalent

(i) $\varepsilon_1$ is \deleted{\illegible{2}} a regular \inserted{section} (i.e. \deleted{\illegible{1}} \inserted{$\ne1$} on each fiber), i.e. $\varepsilon_1=1+N_1$, with $N_1$ a \uline{basis} of $\mathcal N_1$) and $\widetilde\rho^{-1}(x_1)$ is a section of $\unsure{\mathbb P|P}$ \deleted{everywhere} $\ne x_1$ in each fiber.

\margin{\uline{N.B.} We have denoted by $x_1$ the section of $\unsure{\mathbb P|P}$ over $S$ corresponding to $L_1$.}

(ii) \unsure{Ditto|Item}, with $\widetilde\sigma^{-1}$ in place of $\widetilde\rho$.

\margin{$\begin{array}{l}(i\text{ bis})\\(ii\text{ bis})\end{array}\ \big|$ \unsure{Ditto|Item}, with $\widetilde\sigma,\widetilde\rho$ in place of $\widetilde\sigma^{-1},\widetilde\rho^{-1}$.}

(iii) We have (42), i.e. $c=\lambda_0\mathrm{Tr}\,\rho+\mathrm{Tr}\,\sigma$ is \unsure{invertible}.

\uline{Proof.} We have, thanks to the fact that $\widetilde\rho\widetilde\sigma^{-1}\in\Gamma\widetilde L_1$, normalizes $\widetilde L_1$, \inserted{\unsure{or what amounts to the} same, $L_1$,} the relation
\[
\widetilde\rho^{-1}(L_1)=\widetilde\sigma^{-1}(L_1)
\]
or what amounts to the same, we have
\[
\widetilde\rho^{-1}(x_1)=\deleted{\unsure{\widetilde\rho^{-1}}}\inserted{\widetilde\sigma^{-1}}(x_1), \tag{\unsure{22|24}}
\]
which proves the equivalence of (i) and (ii).

Let us show that $(i)\Longleftrightarrow(iii)$. We may assume that $S$ is the spectrum of an algebraically closed field, if it suits us, and we may assume $\lambda_0=1$ (at the cost \deleted{\illegible{1}} of replacing $\rho$ by $\lambda_0\rho$). If we had $\varepsilon_1=1$ i.e. \deleted{$\rho+\sigma$} $1=-\rho\sigma^{-1}$ i.e. $\rho+\sigma=0$, we would have \deleted{$\unsure{\mathrm{Tr}}$} $\mathrm{Tr}(\rho+\sigma)=\mathrm{Tr}\,\rho+\mathrm{Tr}\,\sigma=0$, so (iii) implies the regularity of $\varepsilon_1$ i.e. of $N_1$, \deleted{\illegible{1}} \inserted{we may assume} therefore that $\varepsilon_1$ is regular. We have to show that if we have such a regular unipotent \deleted{\unsure{operator}} \inserted{\unsure{endomorphism}} $\varepsilon_1$\inserted{$=1+N_1$}, of $F$, associated with a $D\subset F$,
\mspage{296}

\scanpage{298}
\mspage{450}

and \unsure{an} \deleted{\unsure{operator} \illegible{1}} \inserted{\unsure{endomorphism}} $\sigma$ of $F$, whence
$\rho=-\varepsilon_1\sigma$, then \deleted{\illegible{2}}
$c=\operatorname{Tr}\sigma+\operatorname{Tr}\rho=\operatorname{Tr}\sigma-\operatorname{Tr}\varepsilon_1\sigma
=-\operatorname{Tr}N_1\sigma$ is invertible iff $D+\sigma^{-1}.D=F$
i.e. $\sigma D+D=F$, \deleted{\unsure{iff}} Let us take a basis
$e_0,e_1$ of $F$ such that $e_1$ generates $D$, and $N_1$ is written as the matrix
$\begin{pmatrix}0&0\\1&0\end{pmatrix}$,
let $\sigma=\begin{pmatrix}a&b\\c&d\end{pmatrix}$ we have
$N_1\sigma=\begin{pmatrix}0&0\\a&b\end{pmatrix}$,
$\operatorname{Tr}N_1\sigma=b$,
$\sigma e_{\unsure{1}}=a e_0+b e_1$, and
$D+\sigma D=D+b\mathcal O_S e_{\unsure{2|0}}$
is equal to $F=\mathcal O_S^2$ iff $b$ invertible, q.e.d.

Let us suppose these equivalent conditions satisfied, \deleted{\unsure{in} $\unsure{P}$} which are also expressed by saying that $L_1$ and
\[
L_0=\widetilde\rho^{-1}(L_1)=\unsure{\widetilde\sigma}^{-1}(L_1)
\tag{23}
\]
are distinct on each fiber, i.e. \unsure{define} Borels $B_1,B_0\subset G$ such that
\[
T=B_1\cap B_0
\tag{24}
\]
is a torus (maximal) of $G$, \deleted{\unsure{which}}
The latter is also the \deleted{\illegible{1}} intersection of the stabilizers $B_0,B_1$ in $G$ of the \inserted{sections}
$x_0,x_1$ of $\unsure{\mathcal P|P}$ which correspond to $L_0,L_1$
(satisfying
\[
x_1=\widetilde\rho x_0=\widetilde\sigma x_0
\tag{25}
\]
), and is \unsure{the inverse} \unsure{image} \unsure{of the subgroup|of the subgroups} (\unsure{privileged|privileged} \inserted{\unsure{toric|toric}}

\mspage{297}

\scanpage{299}

\unsure{called|whose} $\widetilde T$ in $PG=\underline{\operatorname{Aut}}_S(\unsure{\mathcal P|P})$.
We can also say that the datum of two distinct sections \deleted{\unsure{points}}
$x_0,x_1$ of the projective bundle $\unsure{\mathcal P|P}$ \deleted{\unsure{makes it possible to}} amounts to that of the Borel
$B_1$ corresponding to the choice of $x_1$), \uline{and} of the maximal torus $T$ of $B_1$
(which is then written
\[
B_1=T.L_1 \qquad (\unsure{produit}\ \tfrac12\ \text{direct})
\tag{26}
\]

\margin{the pair $(B_1,T)$ is called a \uline{\unsure{Killing} pair}.}

which allows us \unsure{then} to recover $x_0$ as \unsure{one}
\inserted{corresponding to} \unsure{the unique} Borel $B_0$ such that
$T=B_0\cap B_1$.
(The Borels $B_0$ such that $B_0\ne B_1$ on each fiber,
\inserted{i.e. such that $B_0\cap B_1$ is a torus,}
correspond one-to-one to the maximal tori $T$ of $B_1$, by the
\unsure{preceding} formula $T=B_0\cap B_1$.)

This also amounts to: the datum, in addition to $B_1$, \inserted{i.e. to $x_1$,}
of a section \inserted{$x_0$,} of the affine-line bundle
\inserted{$U_1=$}$\unsure{\mathcal P|P}\smallsetminus\{x_1\}$,
\deleted{which \unsure{then makes it possible} to identify $U_1$}
\inserted{i.e. of an \unsure{isomorphism} from $U_1$ onto the}
vector bundle $L_1$ of its translations
\[
U_1\overset{\mathrm{dfn}}{=}\unsure{\mathcal P|P}\smallsetminus\{x_1\}\xrightarrow{\sim}L_1.
\tag{27}
\]

To finish ``pinning'' the situation (i.e.

\mspage{298}

\scanpage{300}
\mspage{451}

reducing $\unsure{\mathcal P|P}$ to the standard projective bundle $\mathbb P^1_S$)
we still have to give ourselves a basis of $L_1$ i.e. a section $x_\infty$ of
\[
L_1^*\simeq U_1\smallsetminus\{x_0\}\simeq\unsure{\mathcal P|P}\smallsetminus\{x_1,x_0\}.
\tag{28}
\]

Now the section $\varepsilon_1$ of $L_1$ is precisely a section of $L_1^*$ -- so we have a complete pinning of the situation.
In terms of $(F,D=D_1)$ corresponding to $L_1$,

\deleted{we have
\[
F=D_1\oplus\underbrace{\sigma^{-1}(D_1)}_{D_0}
=D_1\oplus\underbrace{\rho^{-1}(D_1)}_{D_0},
\]
\unsure{we may assume} (at the cost of tensoring with $D_1^{\otimes-1}$) that $D_1\simeq\mathcal O_S$,
\inserted{\unsure{we have} $N_1|D_0:$}
whence \illegible{1}
\[
\left\{
\begin{aligned}
F&=\mathcal O_S^2\unsure{x}=\mathcal O_S e_1+\mathcal O_S e_0=\mathcal O_S e_0+\mathcal O_S e_1.\\
D_1&=\mathcal O_S e_1,\quad D_0=\mathcal O_S e_0\\
&\unsure{OPS}
\end{aligned}\right.
\tag{\unsure{29}}
\]
we will then have $N_1=(\,\unsure{1}$}
we will have
\[
\left\{
\begin{aligned}
\rho^{-1}(D_1)&=\unsure{\sigma}^{-1}(D_1)\overset{\mathrm{dfn}}{=}D_0\\
F&=D_0\oplus D_1
\end{aligned}\right.
\tag{29}
\]
\deleted{(\unsure{30})}
\[
N_1|D_0:\quad D_0\xrightarrow{\sim}D_1
\tag{30}
\]
and (at the cost of tensoring with $D_1^{\otimes-1}$) \unsure{we may assume}
$D_1\simeq\mathcal O_S$, whence
\[
\left\{
\begin{aligned}
F&\simeq\mathcal O_S^2=\mathcal O_S e_0+\mathcal O_S e_1\\
D_0&=\mathcal O_S e_0,\quad D_1=\mathcal O_S e_1,\\
N_1e_1&=0,\quad N_1e_0=e_1,
\quad\deleted{\text{i.e. }N_1=\begin{pmatrix}0&0\\1&0\end{pmatrix}}
\end{aligned}\right.
\tag{31}
\]

\mspage{299}

\scanpage{301}

i.e.
\[
N_1=\begin{pmatrix}0&0\\1&0\end{pmatrix},\qquad
\varepsilon_1=\begin{pmatrix}1&0\\1&1\end{pmatrix}.
\tag{32}
\]

On the other hand \deleted{\unsure{we will have} $\rho D_0=\sigma\unsure{D}$}
\[
\rho D_0\unsure{=|\ne}\sigma D_0=D_1
\]
\deleted{which is expressed} \inserted{whence}
\[
\rho\sigma^{-1}(D_1)=D_1
\]
\deleted{and we suppose $\rho,\sigma$ chosen over
$\widetilde\rho,\widetilde\sigma$ in such a way that we have
$\rho\sigma^{-1}|D_1=-1$, which (taking into account that
$\widetilde\rho\widetilde\sigma^{-1}=\widetilde\varepsilon_1$)
implies that $\rho\sigma^{-1}(F/D_1)=-1$, i.e.}
\inserted{We \unsure{choose} $\rho,\sigma$ sections in $\mathrm{Gl}(2)$)
over $\widetilde\rho,\widetilde\sigma$
(sections of $\mathrm{GP}(1)\simeq\mathrm{Gl}(2)/\mathrm G_m$
\unsure{satisfying} $\widetilde\rho\widetilde\sigma^{-1}=\widetilde\varepsilon_1$,}
in such a way as to have
\[
\boxed{\rho\sigma^{-1}=-\varepsilon_1}
\tag{33}
\]
which determines each of the two sections $\rho,\sigma$ as a function of the other
(\unsure{=|itself} determined modulo a scalar,
in \deleted{function} \inserted{\unsure{terms}}
of $\unsure{\widetilde\sigma}$ resp. \inserted{of} $\widetilde\rho$).
As
\[
\rho D_0=D_1,
\]
we will have
\[
\rho=\begin{pmatrix}0&\unsure{u|U}\\\unsure{v|V}&\unsure{w|W}\end{pmatrix},
\qquad \det\rho=-\unsure{uv|UV}\ \text{invertible}
\tag{*}
\]
and we can remove the scalar \unsure{ambiguity|ambiguity} in $\rho$
by \deleted{\unsure{writing}} normalizing by
\[
\rho=\begin{pmatrix}0&1\\\unsure{v|V}&\unsure{w|W}\end{pmatrix}
\qquad \unsure{v|V}\in\Gamma\mathrm G_m.
\tag{34}
\]

\mspage{300}
